\section{Introduction}
\label{sec:intro}

Surrogate-assisted optimization has become an important strategy for accelerating structural design problems that require repeated finite-element analysis (FEA). In shear wall structural design for high-rise buildings, engineers need to determine wall thicknesses, beam section dimensions, material grades, and other design parameters while satisfying multiple code-based constraints, including inter-story drift, torsional response, axial compression ratio, shear-compression ratio, and member reinforcement requirements. For each candidate design, reliable evaluation usually requires parametric modeling, FEA, structural response extraction, and member-level code checking. When population-based optimizers such as genetic algorithms or particle swarm optimization are adopted, a large number of candidate designs must be evaluated, making repeated FEA and code checking the dominant computational bottleneck.

Existing studies have explored various approaches to improve the automation and efficiency of shear wall structural design. Traditional optimization methods, such as genetic algorithms, tabu search, and other metaheuristic algorithms, have been used to search for cost-effective shear wall layouts or section configurations under engineering constraints . More recent studies have introduced prior design knowledge, CAD/BIM-based automatic modeling, machine learning, and generative models into preliminary shear wall design, enabling more efficient layout generation, structural scheme exploration, and data-driven design support . These studies have advanced the automation of shear wall structural design, but many of them still rely on expensive structural analyses to evaluate candidate solutions. Therefore, for section and material optimization of a given shear wall layout, the key challenge is not only how to search the design space, but also how to allocate limited FEA evaluations to candidates that are more likely to be feasible and cost-effective.

Surrogate models provide a natural way to reduce the number of expensive evaluations in structural optimization. Response surface models, Kriging models, support vector machines, neural networks, and other data-driven surrogates have been widely used to approximate structural responses, replace costly simulations, or guide optimization and reliability analysis . In building engineering, surrogate models have also been applied to predict seismic responses, wind-induced responses, material quantities, and performance indicators of structural systems. These methods demonstrate the potential of data-driven prediction for accelerating structural design. In shear wall structural optimization, a single static surrogate has difficulty reproducing the coupled global and member-level checks that govern design feasibility.

Several limitations become particularly important when surrogate models are used for shear wall structural optimization. First, the behavior of shear wall buildings is strongly layout-dependent. Even under the same section and material parameters, different wall arrangements may lead to different lateral stiffness distributions, load paths, seismic responses, and reinforcement demands. A global surrogate trained on multiple layouts may therefore suffer from layout-specific prediction bias when applied to a new design case. Second, optimization is an iterative process in which the candidate distribution gradually changes. A surrogate model trained offline may not remain accurate in the local region explored by the optimizer. Third, different surrogate predictions should serve different roles in the optimization process. Feasibility prediction should be conservative to avoid rejecting potentially feasible designs, whereas reinforcement usage prediction is more useful for ranking candidates by material economy. These requirements are not fully addressed by conventional surrogate modeling settings that focus mainly on global classification or regression accuracy.

This paper proposes a calibrated surrogate-assisted optimization framework for FEA-efficient shear wall structural design. The proposed framework uses surrogate models as pre-evaluation tools for candidate screening and analysis budget allocation. A layout-conditioned surrogate model is first trained to estimate candidate feasibility and reinforcement usage from layout information and design parameters. During optimization, FEA-verified samples accumulated from the current layout are further used to perform online local calibration, so that global surrogate predictions can be corrected for the current design case and local search region. The calibrated feasibility and reinforcement predictions are then used to conservatively filter low-value candidates and prioritize promising candidates for true FEA and code checking.

The main contributions of this study are summarized as follows. First, a layout-conditioned surrogate modeling strategy is developed for candidate assessment in shear wall section and material optimization, enabling both feasibility screening and reinforcement usage estimation before FEA. Second, an online local calibration mechanism is introduced to correct layout-specific surrogate bias using FEA-verified samples generated during the optimization process. Third, a calibrated surrogate-assisted optimization framework is established to allocate FEA budgets more effectively through conservative candidate screening and priority ranking.

The remainder of this paper is organized as follows. Section 2 defines the shear wall structural optimization problem and introduces the calibrated surrogate modeling strategy. Section 3 presents the surrogate-assisted optimization framework and the online update mechanism. Section 4 describes the dataset, experimental settings, and evaluation metrics. Section 5 reports and discusses the experimental results. Section 6 concludes the paper and outlines future research directions.